\section{Evaluation Design}\label{sec:evaluation}
Section~\ref{sec:method-system} specified how \textit{Boyun} executes a world;
this section specifies how we test whether that execution does what it claims
and what it adds in practice. Two kinds of test are needed, and they answer
different questions. \textit{Verification} asks whether the program carries
out the stated rules correctly. \textit{Validation} asks whether those rules
and the model's choices resemble real actors and markets. The tests
below are mainly verification and capability tests; none of them validates
\textit{Boyun} as a forecaster, and the results in Section~\ref{sec:results}
are interpreted accordingly.

\subsection{Research Questions}
Given that boundary, the evaluation addresses six research questions (RQs),
ordered from the most controlled test to the most open comparison:
\begin{enumerate}\setlength{\itemsep}{1pt}
\item[\textit{RQ1}] Do the implemented calculations, timing, accounting, and
replay match their specifications in Section~\ref{sec:method-system}?
\item[\textit{RQ2}] Do information delivery, joint rules, and numeric
preferences have separate and inspectable effects on choices and outcomes?
\item[\textit{RQ3}] Can a fourteen-day financial case with sixteen actors run
with settled trades, recorded rejections, matched policy branches, replay,
and a bounded budget, and what does it record?
\item[\textit{RQ4}] Given the same brief, what does \textit{Boyun} record that
the original \textit{MiroFish} does not, and what does \textit{MiroFish}
provide that \textit{Boyun} lacks?
\item[\textit{RQ5}] What do the local model's score distributions and
regime responses reveal about the limits of its simulated behavior?
\item[\textit{RQ6}] In a ninety-day lending launch, which observations follow
from explicit rules and which depend on unvalidated actor choices?
\end{enumerate}

\subsection{Synthetic Mechanism Experiments}\label{sec:eval-synthetic}
The first two questions require tests whose correct answers are known in
advance. Six small synthetic experiments provide them
(Table~\ref{tab:experiments}). Each uses invented inputs, and wherever an
action is needed it uses a fixed policy rather than a model, so any
discrepancy must come from the engine. No model is called and network access
is disabled while they run. Experiments~1, 2, and~4 test the analytical tools
and scoring formulas against hand calculations; Experiment~3 tests timing,
accounting, and replay; Experiments~5 and~6 address RQ2 by switching single
mechanisms on and off. These experiments are complemented by the engine's
automated test suite.

\begin{table}[H]
\caption{Synthetic mechanism experiments.}
\label{tab:experiments}
\centering\small
\renewcommand{\arraystretch}{1.12}
\begin{tabularx}{\columnwidth}{@{}>{\raggedright\arraybackslash}p{0.07\columnwidth}>{\raggedright\arraybackslash}p{0.27\columnwidth}Y@{}}
\toprule
Exp. & Mechanism & Known answer checked \\
\midrule
1 & Game analysis & Expected utility, regret, and pure Nash profiles, Eqs.~\eqref{eq:eu}--\eqref{eq:regret} \\
2 & Bayesian updating & Two sequential posteriors, Eq.~\eqref{eq:bayes} \\
3 & Timing and accounting & Event order, clipping, conservation, replay, Eqs.~\eqref{eq:cash-update}--\eqref{eq:conservation} \\
4 & Forecast scoring & Brier and log scores against a baseline, Eqs.~\eqref{eq:brier}--\eqref{eq:logloss} \\
5 & Communication ablation & Delivery and a joint rule switched on and off \\
6 & Preferences vs.\ prose & Changing a weight versus changing persona text \\
\bottomrule
\end{tabularx}
\end{table}

\subsection{MACRON Case Study}\label{sec:eval-macron}
Synthetic checks cannot show whether the engine copes with a realistic case,
so RQ3 uses \textit{MACRON}, a prospective fourteen-day launch of a
hypothetical political memecoin on a \textit{Solana} launch venue, beginning
2 November 2026. A founder with \$50{,}000 of capital, of which \$10{,}000 is a
protected reserve, launches the token on a bonding curve that later migrates
to a pool (Section~\ref{sec:market}). Sixteen actors take part: the founder, a
market maker, four retail traders with different strategies, a large buyer, an
attention trader, four promoters, a meme curator, an organic supporter, a
campaign operator, and a foundation. Each receives its own state, the public
market history, and only the private instructions it is permitted to see.
Real dated events, such as scheduled parliamentary sessions, enter as
exogenous events, without assuming their outcomes.

The case compares founder plans in a full factorial design. The founder's
initial position is \$5{,}000, \$10{,}000, or \$15{,}000; upfront marketing and
operating costs are \$5{,}000 or \$10{,}000; the market regime is
\textit{BULL} or \textit{BEAR}; and the response to an exit alert is either
\textit{pause-and-observe}, which stops discretionary spending and keeps the
tokens, or \textit{staged exit}, which sells in tranches. The regimes are
hypotheses about risk appetite, not observed market states. The resulting
$3\times2\times2\times2=24$ paths run hourly for up to 336 hours, with finer
steps around the launch. Each path uses one seed, so the design compares plans
under common conditions rather than sampling a population of outcomes. \begin{revised}
All 24 paths use local \textit{Laya} inference from hour 0 to hour 336.
No hosted-model decisions or mixed-model continuations enter these case results.
The packages, seeds, response policies, and evidence registers are fixed before
execution. Offline reference paths remain engineering checks rather than
substitutes for these model-driven paths.
\end{revised}

\subsection{Comparison with \textit{MiroFish}}\label{sec:eval-mirofish}
Because \textit{Boyun} was inspired by \textit{MiroFish}, RQ4 compares the two
on the same case. The original \textit{MiroFish} (source revision
\texttt{7657031}) \cite{mirofish} received the same frozen \textit{MACRON}
brief and event calendar, generated its sixteen personas with its own native
prompts, and ran its social simulation on the \textit{OASIS} platform
\cite{yang2025} for 336 hourly rounds in each regime, followed by its own
report generation. It used \textit{Gemini} 3.5 Flash-Lite on free quotas.
Interruptions caused by quota limits required restoring the saved actor
memories and databases; the original random state could not be recovered and
was reseeded, so the run is a declared recovery-adapted version rather than
an uninterrupted stock run.

The comparison is therefore matched in its inputs but not in everything else.
The two systems use different actor models, budgets, and actor activation
rules, and they are designed to produce different kinds of output, social
posts in one case and settled trades in the other. We therefore compare what
each system records and reports, rather than ranking them on a single score,
and we make no claim that either forecasts the real launch.

\subsection{Local Actor Behavior}\label{sec:eval-laya}
\begin{revised}
Beyond completion, RQ5 examines how local actor scores translate into
observed behavior. We report score concentration, filled and rejected trades,
and differences between the declared \textit{BULL} and \textit{BEAR} regimes.
The lifecycle engine samples actions from the saved score distributions;
therefore a weak preference between options can produce substantial turnover.
These are descriptive diagnostics of this model and selection rule, not a
calibration study or a measure of real investor behavior.
\end{revised}

\subsection{NEAR Lending Case Study}\label{sec:eval-lending}
Apart from token trading, RQ6 changes the financial mechanism to collateralized
lending, retaining local \textit{Laya} inference. Its case
is a prospective launch of a lending market on \textit{NEAR} that accepts
native NEAR as collateral and lends Circle's native USDC, evaluated over its
first 90 days. The world was built with the user through 41 recorded answers.
Because no partner terms had been negotiated, the user chose to assume that
\$1{,}000{,}000 of USDC from friendly partners is already deposited at
opening. The team holds \$100{,}000, of which \$10{,}000 is a capped loss
reserve, and suppliers earn points only on USDC that is actually borrowed.
Twenty decision-makers take part: the team, four friendly partners, two large
independent suppliers and two supplier groups, four borrowers and one
borrower group, an actor who both supplies and borrows to earn points,
two liquidators, and three distribution partners. Their balances are
synthetic, sized against the user's premises and a dated snapshot of
\textit{RHEA}, the main existing NEAR lending venue, which paid 9.88\% a year to
native-USDC suppliers and charged borrowers 15.29\% \cite{rhealending}. Every
actor sees these rates as a fixed outside option.

Within that cast, the rules of Section~\ref{sec:lending} apply with a 50\%
maximum loan-to-value, liquidation at 65\%, a 5\% liquidation bonus, a
\$500{,}000 aggregate borrowing cap, and, in the baseline, a \$50{,}000 cap per
borrower. Outsiders can act only after a channel they follow delivers a
campaign item; paid partners may accept, delay, or decline, and cash is spent
only on delivery, so no conversion rate is assumed. The ledger settles hourly,
and actors who are aware of the launch decide daily and whenever a trigger
fires, such as a NEAR price move of 8\% within 24 hours. Prices are real
hourly closes from Binance \cite{binanceklines}, rescaled so that each window
opens at the snapshot price: a choppy window beginning 12 August 2025, which
contains the market-wide fall of 10 October 2025, and a bear window beginning
9 November 2025.

\begin{revised}
These inputs define four local \textit{Laya} worlds
(Table~\ref{tab:lending-design}). Codes identify the price path (C for choppy,
B for bear), per-borrower cap in thousands of dollars, rate curve, and local
model. Each world has one seed and there is no no-launch arm; real-world odds
are unknown. Raising the cap also splits a borrower cohort, so that comparison
changes two design factors and does not isolate the cap alone.
\end{revised}

\begin{table}[t]
\caption{\rev{Local \textit{Laya} lending worlds.} R2: 2\% base rate, 10\% at 80\% utilization,
50\% at full use. R8: 8\% base, 15\% at 50\%, 60\% at full use. Worlds with a
\$100K cap also split the 80-person borrower group into eight accounts.}
\label{tab:lending-design}
\centering\footnotesize\color{revisedred}
\renewcommand{\arraystretch}{1.1}
\begin{tabular}{@{}lllll@{}}
\toprule
World & Path & Cap & Curve & Model \\
\midrule
C-P50-R2-L & Choppy & \$50K & R2 & \textit{Laya} \\
C-P100-R2-L & Choppy & \$100K & R2 & \textit{Laya} \\
C-P100-R8-L & Choppy & \$100K & R8 & \textit{Laya} \\
B-P100-R8-L & Bear & \$100K & R8 & \textit{Laya} \\
\bottomrule
\end{tabular}
\end{table}

\subsection{Metrics and Reporting}\label{sec:eval-metrics}
The metrics follow from the questions and are reported with their units and
coverage. Mechanism tests report exact values against hand calculations.
Case runs report fills, rejections, traded cash, local model calls, recorded
run time, and API fees separately from hardware and electricity costs. Because live paths may stop at different hours, coverage is
reported as completed path-hours out of the planned $24\times336$, and every
path is labeled completed, budget-limited, or interrupted. Policies are
compared only at a common time. For each plan pair $p$, let $\tau_p$ be the
latest hour saved in both of its branches, and let $C_p(\tau_p)$ be the
founder's realized net-cash position relative to initial capital at that hour. The matched difference is
\begin{equation}
D_p=C^{\mathrm{exit}}_p(\tau_p)-C^{\mathrm{pause}}_p(\tau_p).
\label{eq:matched}
\end{equation}
A positive $D_p$ means staged exit held more cash than pausing at the same
moment. It counts only realized cash; unsold tokens are reported separately,
so $D_p$ is not a comparison of total wealth.

Where probabilities can be compared with outcomes, \textit{Boyun} uses the
Brier and logarithmic scores \cite{brier1950,gneiting2007}. For $n$ binary
cases with forecast probability $p_j$ and outcome $y_j\in\{0,1\}$,
\begin{align}
\operatorname{Brier}&=\frac{1}{n}\sum_{j=1}^{n}(p_j-y_j)^2,\label{eq:brier}\\
\operatorname{LogLoss}&=-\frac{1}{n}\sum_{j=1}^{n}\big[y_j\log p_j+(1-y_j)\log(1-p_j)\big].\label{eq:logloss}
\end{align}
Both are averages, lower is better, and both are minimized in expectation by
reporting the true probability; the logarithmic score penalizes confident
errors much more heavily. To avoid $\log 0$, probabilities are bounded to
$[10^{-15},1-10^{-15}]$. Each case path has one seed, so we report observed
ranges rather than confidence intervals.
For the lending case, the measures are borrowed principal, utilization,
liquidations, bad debt, the protocol's share of interest against campaign
spending, and how often \textit{Laya} borrowers chose to borrow when offered the option.
